\section{Avaliação da família WiSARD}
\label{sec:wisard}

Esta seção apresenta o \emph{sweep} de hiperparâmetros dos cinco
modelos WiSARD-family (WiSARD, ClusWiSARD, BloomWiSARD, BTHOWeN e
ULEEN) sob o mesmo protocolo \emph{paper-faithful} dos
\emph{baselines}. A varredura, com tetos de capacidade
\emph{unificados} entre os modelos para uma comparação justa
(\S{}\ref{sec:wisard_fair}), avaliou 30\,020 configurações
(30\,000 \emph{per-device} mais 20 na tarefa combinada), cada uma sob
validação cruzada de 4 \emph{folds} e um conjunto \emph{held-out} de
20\% reservado por \texttt{StratifiedShuffleSplit} (semente 42),
totalizando cerca de 150\,mil \emph{fits}. Os resultados brutos estão
em \texttt{results/\_consolidado/grid/wisard\_grid\_*.jsonl}
(30\,020 linhas), com F1 \emph{weighted} de validação cruzada (média e
desvio sobre os 4 \emph{folds}) e do \emph{held-out} em cada linha.

As subseções percorrem os \emph{grids} varridos por modelo, o melhor F1
por (modelo, base) e a generalização no \emph{held-out}; em seguida a
sensibilidade aos termômetros, o \emph{deep dive} de hiperparâmetros e o
\emph{ranking} por base; e encerram com o custo do \emph{sweep} e com a
unificação dos tetos de capacidade (\S{}\ref{sec:wisard_fair}) que torna
a comparação intra-família justa.

\subsection{Hiperparâmetros varridos}

A Tabela~\ref{tab:wisard_grids} resume os \emph{grids} de cada
modelo. Os três modelos baseados em termômetro (WiSARD, ClusWiSARD,
BloomWiSARD) compartilham o mesmo \emph{front-end} de binarização
(termômetro $\times$ \emph{bits} por \emph{feature}), enquanto
BTHOWeN e ULEEN usam seus codificadores internos. O termômetro varia
sobre um conjunto curado de seis variantes \emph{não supervisionadas}:
\texttt{\{linear, distributive, gaussian, exponential, logarithmic,
non-uniform\}}. O \texttt{non-uniform} é um \texttt{DynamicThermometer}
cujo orçamento de \emph{bits} é alocado por \emph{feature} de forma
proporcional à dispersão (desvio-padrão), com piso de 2 \emph{bits} e
orçamento total conservado. Descartamos o \texttt{stochastic} (cujo
\texttt{optimize()} nunca era acionado no \emph{pipeline}, tornando-o
idêntico ao \texttt{distributive}) e o \texttt{supervised} (que usa os
rótulos e quebraria a comparação \emph{unsupervised}).

\begin{table}[ht]
  \centering
  \caption{Hiperparâmetros varridos por modelo (\emph{grid} estendido
           com tetos unificados). Os três primeiros compartilham
           termômetro $\in$ \texttt{\{linear, distributive, gaussian,
           exponential, logarithmic, non-uniform\}} e
           \emph{bits/feature} $\in \{8, 16, 32, 48, 64\}$.}
  \label{tab:wisard_grids}
  \small
  \begin{tabular}{p{0.18\linewidth} p{0.74\linewidth}}
    \toprule
    \textbf{Modelo} & \textbf{Hiperparâmetros} \\
    \midrule
    WiSARD       & \emph{address size} $\in \{2,4,8,16,32,48,64\}$, \texttt{ignore\_zero} $\in \{$T,F$\}$ \\
    ClusWiSARD   & \emph{address size} $\in \{4,8,16,24,32,48,64\}$, \texttt{min\_score} $\in \{0{,}1; 0{,}2\}$, \texttt{threshold} $\in \{10,100\}$, \texttt{limit} $\in \{5,20\}$ \\
    BloomWiSARD  & \emph{address size} $\in \{8,16,24,32,48,64\}$, \texttt{filter\_size} $\in \{256,1024,4096\}$, \texttt{num\_hashes} $\in \{2,3,4\}$ \\
    BTHOWeN      & \emph{address size} $\in \{6,12,16,20,24\}$, \texttt{num\_bits} $\in \{128,256,512,1024,2048,4096\}$, \texttt{num\_hashes} $\in \{2,4,6\}$, \texttt{bits\_per\_input} $\in \{4,8,16,24,32,48\}$ \\
    ULEEN        & \texttt{bits\_per\_input} $\in \{4,8,16,24,32\}$, \texttt{filter\_inputs} $\in \{3,6,12\}$, \texttt{filter\_entries} $\in \{16,64,128,256\}$, \texttt{n\_submodels} $\in \{1,3,5\}$, \emph{epochs} $\in \{5,10,20\}$, \emph{lr} $\in \{0{,}001; 0{,}01\}$ \\
    \bottomrule
  \end{tabular}
\end{table}

\subsection{Best F1 por (modelo, base)}

A Tabela~\ref{tab:wisard_best_f1} reporta o melhor F1 \emph{weighted}
médio (validação cruzada de 4 \emph{folds}) de cada modelo em cada
base, sob a tarefa binária \emph{per-device}. Para cada par (modelo,
base), o ``melhor'' significa a configuração de hiperparâmetros
vencedora dentro do \emph{grid}.

\begin{table}[ht]
  \centering
  \caption{Best F1 \emph{weighted} médio (4 \emph{folds}) por (modelo,
           base), binário \emph{per-device}. Negrito = melhor
           WiSARD da base; última coluna = média sobre as 7 bases.}
  \label{tab:wisard_best_f1}
  \scriptsize
  \begin{tabular}{lcccccccc}
    \toprule
    \textbf{Modelo} & \textbf{Fridge} & \textbf{Garage\_D.} & \textbf{GPS\_T.} & \textbf{Modbus} & \textbf{Motion\_L.} & \textbf{Thermo.} & \textbf{Weather} & \textbf{Média} \\
    \midrule
    WiSARD       & \textbf{1{,}000} & \textbf{1{,}000} & 0{,}876 & 0{,}953 & 0{,}503 & 0{,}505 & 0{,}932 & \textbf{0{,}824} \\
    ClusWiSARD   & \textbf{1{,}000} & \textbf{1{,}000} & \textbf{0{,}878} & \textbf{0{,}954} & 0{,}475 & 0{,}511 & \textbf{0{,}934} & 0{,}822 \\
    BloomWiSARD  & \textbf{1{,}000} & \textbf{1{,}000} & 0{,}877 & 0{,}637 & 0{,}475 & 0{,}508 & 0{,}852 & 0{,}764 \\
    BTHOWeN      & \textbf{1{,}000} & \textbf{1{,}000} & 0{,}868 & 0{,}590 & \textbf{0{,}531} & 0{,}511 & 0{,}818 & 0{,}760 \\
    ULEEN        & \textbf{1{,}000} & \textbf{1{,}000} & 0{,}832 & 0{,}614 & \textbf{0{,}531} & \textbf{0{,}512} & 0{,}816 & 0{,}758 \\
    \bottomrule
  \end{tabular}
\end{table}

Observações:
\begin{itemize}[leftmargin=*,itemsep=0pt]
  \item \texttt{Fridge} e \texttt{Garage\_Door}: \emph{todos} os
        cinco modelos atingem 1{,}000, mesma saturação observada
        nos \emph{baselines}. As \emph{features} de \emph{leak}
        (\texttt{temp\_condition}, \texttt{sphone\_signal}) são
        capturadas com facilidade pela discretização do termômetro.
  \item \texttt{GPS\_Tracker}, \texttt{Weather}: a família converge
        para 0{,}83--0{,}88 (GPS) e 0{,}82--0{,}93 (Weather); com os
        tetos de endereço estendidos, WiSARD e ClusWiSARD em
        \texttt{Weather} (0{,}93) aproximam-se dos \emph{baselines}
        (\textbf{kNN} = 0{,}916, \textbf{RF} = 0{,}975).
  \item \texttt{Modbus}: WiSARD e ClusWiSARD agora \emph{empatam} no
        topo (0{,}953/0{,}954, ambos com \emph{address size} = 64 e 64
        \emph{bits}; nessa capacidade máxima \texttt{linear} e
        \texttt{distributive} ficam estatisticamente indistinguíveis,
        $\Delta < 0{,}001 \ll$ desvio entre \emph{folds}), bem acima
        de BloomWiSARD (0{,}64), ULEEN (0{,}61) e BTHOWeN (0{,}59). O
        endereçamento amplo recupera a estrutura não-linear dos códigos
        FC e \emph{alcança ou supera} os modelos de árvore
        (RF/CART = 0{,}94), ganho que a varredura original perdia por
        limitar \emph{address size} a 32 (WiSARD) e 16 (ClusWiSARD).
  \item \texttt{Thermostat}: dificuldade compartilhada com os
        \emph{baselines} (0{,}50--0{,}51 para todos os modelos).
  \item \texttt{Motion\_Light}: colapso para predição majoritária
        (0{,}475--0{,}531) em todos os modelos: BTHOWeN e ULEEN são
        os únicos a romper marginalmente o platô (0{,}531).
\end{itemize}

WiSARD (0{,}824) e ClusWiSARD (0{,}822) lideram a média
\emph{per-device} \emph{praticamente empatadas}, seguidas de
BloomWiSARD (0{,}764), BTHOWeN (0{,}760) e ULEEN (0{,}758). Essa
paridade é resultado da unificação dos tetos de capacidade
(\S{}\ref{sec:wisard_fair}): ClusWiSARD vence \texttt{Modbus},
\texttt{Weather} e \texttt{GPS\_Tracker}, e só não lidera a média
porque WiSARD não colapsa para predição majoritária em
\texttt{Motion\_Light} (0{,}503 vs.\ 0{,}475). Na varredura original,
com tetos de endereço desiguais (32 para WiSARD, 16 para ClusWiSARD),
ClusWiSARD aparecia $0{,}04$ atrás, um artefato de cobertura, não de
modelo.

\subsection{Generalização no \emph{held-out}}

A Tabela~\ref{tab:wisard_holdout} compara a média \emph{per-device} de
F1 \emph{weighted} entre a validação cruzada e o conjunto
\emph{held-out} de 20\%, nunca visto durante a seleção de
hiperparâmetros. Para WiSARD, ClusWiSARD e BloomWiSARD as duas medidas
praticamente coincidem (gap $\le$ 0{,}002), indício de seleção sem
\emph{overfitting}. BTHOWeN é o que mais cai
(0{,}760 $\rightarrow$ 0{,}725), seguido de ULEEN
(0{,}758 $\rightarrow$ 0{,}739), concentrado nas bases estruturadas,
onde o ajuste do \emph{bleaching} por busca ternária interna (no
\emph{fit}, não no \emph{grid}) superajusta aos \emph{folds}.

\begin{table}[ht]
  \centering
  \caption{Média \emph{per-device} (7 bases) de F1 \emph{weighted}:
           validação cruzada vs.\ \emph{held-out} 20\%.}
  \label{tab:wisard_holdout}
  \small
  \begin{tabular}{lccccc}
    \toprule
    & \textbf{WiSARD} & \textbf{ClusW.} & \textbf{BloomW.} & \textbf{BTHOWeN} & \textbf{ULEEN} \\
    \midrule
    Validação cruzada & 0{,}824 & 0{,}822 & 0{,}764 & 0{,}760 & 0{,}758 \\
    \emph{Held-out} 20\% & 0{,}825 & 0{,}823 & 0{,}763 & 0{,}725 & 0{,}739 \\
    \bottomrule
  \end{tabular}
\end{table}

\subsection{Sensibilidade aos termômetros}

A Figura~\ref{fig:thermo_boxplot} reporta a distribuição de F1 ao
variar o termômetro, agregada sobre as 7 bases e todas as
configurações dos três modelos baseados em termômetro. O
\texttt{distributive} domina a distribuição agregada de F1 nos três
modelos, com o \texttt{gaussian} em segundo. Sua vantagem sobre o
\texttt{linear} (cortes equiespaçados) é maior em distribuições de cauda
longa e sob orçamento de \emph{bits} escasso: em \texttt{Weather} o
\texttt{distributive} lidera por $\approx 0{,}06$ em \emph{todos} os
orçamentos, inclusive a 64 \emph{bits} (0{,}932 contra 0{,}872), porque
a binarização equiespaçada subdivide pouco a região de massa.
\texttt{Modbus} é a exceção reveladora: lá o \texttt{distributive}
supera o \texttt{linear} em $\approx 0{,}09$ a 8 \emph{bits}, mas a
diferença encolhe com o orçamento e some a 64 \emph{bits} (ambos
$\approx 0{,}95$), quando já há resolução para separar os códigos FC
mesmo com cortes uniformes. O \texttt{non-uniform} e o \texttt{exponential}
ficam no fim do \emph{ranking}: alocar \emph{bits} por dispersão não
compensa quando poucas \emph{features} concentram a variância útil.

\begin{figure}[ht]
  \centering
  \includegraphics[width=0.85\linewidth]{v06_thermo_boxplot.png}
  \caption{Distribuição de F1 \emph{weighted} por termômetro,
           agregada sobre as 7 bases e todas as configurações dos três
           modelos baseados em termômetro. \texttt{distributive} e
           \texttt{gaussian} dominam.}
  \label{fig:thermo_boxplot}
\end{figure}

\subsection{Deep dive por modelo}

As Figuras~\ref{fig:dd_wisard} a~\ref{fig:dd_uleen} mostram, para cada
modelo, o \emph{landscape} do \emph{sweep} (F1 médio sobre bases ao
variar os dois hiperparâmetros principais) e a melhor configuração por
base. O padrão é o mesmo nos cinco modelos: o eixo de capacidade domina
o \emph{landscape} e o F1 cresce de forma quase monotônica ao longo
dele (\emph{address size} no trio de termômetro, \texttt{bits\_per\_input}
em BTHOWeN e ULEEN), com a melhor configuração por base encostando no
teto (64 para o trio; 48 e 32 \emph{bits} para BTHOWeN e ULEEN). Os
demais eixos (termômetro, \texttt{num\_hashes}, \texttt{filter\_inputs},
\emph{lr}) são de segunda ordem. É esse comportamento de memorizador,
visível em todos os painéis, que motiva a unificação de tetos da
\S{}\ref{sec:wisard_fair}.

\begin{figure}[!htb]
  \centering
  \begin{subfigure}[t]{\linewidth}
    \includegraphics[width=\linewidth]{v06_dd_wisard.png}
    \caption{WiSARD.}
    \label{fig:dd_wisard}
  \end{subfigure}\\[0.6em]
  \begin{subfigure}[t]{\linewidth}
    \includegraphics[width=\linewidth]{v06_dd_cluswisard.png}
    \caption{ClusWiSARD.}
  \end{subfigure}\\[0.6em]
  \begin{subfigure}[t]{\linewidth}
    \includegraphics[width=\linewidth]{v06_dd_bloomwisard.png}
    \caption{BloomWiSARD.}
  \end{subfigure}
  \caption{Deep dive de hiperparâmetros dos modelos WiSARD-family
           (parte 1/2). Cada painel mostra o \emph{landscape} do
           \emph{sweep} (F1 médio em função dos dois hiperparâmetros
           principais), a melhor configuração por base e a
           distribuição de F1 entre configs.}
\end{figure}

\begin{figure}[!htb]
  \centering
  \begin{subfigure}[t]{\linewidth}
    \includegraphics[width=\linewidth]{v06_dd_bthowen.png}
    \caption{BTHOWeN.}
  \end{subfigure}\\[0.6em]
  \begin{subfigure}[t]{\linewidth}
    \includegraphics[width=\linewidth]{v06_dd_uleen.png}
    \caption{ULEEN.}
    \label{fig:dd_uleen}
  \end{subfigure}
  \caption{Deep dive de hiperparâmetros de BTHOWeN e ULEEN (parte 2/2).
           O F1 cresce com a capacidade (\texttt{bits\_per\_input}) e a
           melhor configuração por base senta no teto, mesmo padrão do
           trio de termômetro.}
\end{figure}

\subsection{Ranking dos modelos por base}

A Figura~\ref{fig:ranking_per_base} põe os cinco modelos lado a lado
por base (altura = best F1, anotação = tempo de \emph{fit} único). As
alturas repetem a Tabela~\ref{tab:wisard_best_f1}; o que a figura
acrescenta é o eixo de custo, que separa os modelos em duas ordens de
grandeza: BTHOWeN treina mais rápido ($\approx 0{,}09$--$0{,}10$\,s), o
trio em C++ (WiSARD, ClusWiSARD, BloomWiSARD) fica em
$0{,}4$--$0{,}8$\,s, e o ULEEN é o \emph{outlier} a $6{,}4$--$9{,}0$\,s.
A \S{}\ref{sec:wisard_custo} mostra que essa lentidão é de
\emph{implementação} (referência em Python vs.\ C++), não de algoritmo.

\begin{figure}[ht]
  \centering
  \includegraphics[width=0.92\linewidth]{v06_ranking_per_base.png}
  \caption{Ranking dos cinco modelos WiSARD-family por base
           (best F1, anotação de tempo médio de \emph{fit}).}
  \label{fig:ranking_per_base}
\end{figure}

\subsection{Custo e protocolo do \emph{sweep}}
\label{sec:wisard_custo}

A varredura completa (30\,020 configurações) roda em cerca de 6--7\,h
de \emph{wall-clock} nesse \emph{hardware}, com as sete bases
\emph{per-device} executadas em paralelo. O custo é dominado pelo
ULEEN, cuja implementação de referência em \emph{Python} é cerca de
uma ordem de magnitude (13--16$\times$) mais lenta no treino que o trio
WiSARD/ClusWiSARD/BloomWiSARD em C++, uma ressalva de
\emph{implementação}, não de algoritmo, retomada na discussão de custo
(\S{}\ref{sec:comparacao}).
Três decisões de engenharia tornaram esse orçamento viável e
reprodutível. Primeiro, o \emph{cache} de codificação: cada
\emph{encoding} (termômetro, \emph{bits}, \emph{fold}) é computado uma
única vez e reutilizado por todas as configurações internas dos três
modelos baseados em termômetro. Segundo, a poda de regiões explosivas
do \emph{grid} (ex.: ClusWiSARD com \emph{address size} e
\texttt{limit} simultaneamente altos) e a validação de
\emph{address size} contra o número de \emph{bits} disponível, que
explica a leve variação de cobertura entre bases com mais
\emph{features}. Terceiro, o conjunto \emph{held-out} de 20\% reservado
antes de qualquer seleção, reportado na Tabela~\ref{tab:wisard_holdout}
como verificação independente de generalização.

\subsection{Tetos de capacidade unificados e \emph{tuning} justo}
\label{sec:wisard_fair}

Uma primeira versão do \emph{sweep} usava tetos de \emph{address size}
\emph{desiguais} entre os modelos (32 para WiSARD, 24 para BloomWiSARD,
apenas 16 para ClusWiSARD) e \emph{bits/feature} no máximo 32. Como as
redes sem peso são memorizadoras (mais capacidade de endereçamento
quase sempre adiciona sinal nas bases estruturadas), a melhor
configuração de cada modelo \emph{encostava no teto}, e os tetos
desiguais penalizavam injustamente o ClusWiSARD: em \texttt{Modbus} ele
ficava em 0{,}603 contra 0{,}817 da WiSARD, uma diferença que parecia
de \emph{modelo} mas era de \emph{cobertura de grid}.

Re-executamos a varredura com tetos \textbf{unificados} em 64
(\emph{address size} e \emph{bits}) para todo o trio de termômetro, e
ampliamos analogamente os eixos de BTHOWeN e ULEEN
(Tabela~\ref{tab:wisard_grids}). O efeito é direto: o ClusWiSARD em
\texttt{Modbus} sobe para 0{,}954 (empatando com a WiSARD em 0{,}953),
e a média \emph{per-device} dos dois fica em $0{,}822$ vs.\ $0{,}824$,
um empate técnico onde antes havia $0{,}04$ de diferença.

A ressalva honesta é que, nas bases estruturadas, a melhor
configuração ainda fica na fronteira (64/64): para memorizadores o
ótimo de validação cruzada tende a sentar-se no teto de capacidade, e
``estender até o interior'' não termina. Uma sondagem em 128 quantifica
o retorno decrescente: o salto 32$\rightarrow$64 rendeu $+0{,}135$ de
F1 em \texttt{Modbus}, enquanto 64$\rightarrow$128 rende apenas
$+0{,}011$ (\texttt{Modbus}, já convergido) a $+0{,}027$
(\texttt{Weather}), com ganho idêntico para WiSARD e ClusWiSARD,
isto é, abaixo de $0{,}03$ e sem alterar o \emph{ranking}. Reportamos
portanto os resultados no orçamento generoso e \emph{consistente} de
64: como o teto é o mesmo para todos os modelos do trio, a comparação
intra-família é \emph{não enviesada}, ainda que o valor absoluto possa
subir poucos centésimos com mais capacidade.
